\documentclass[11pt,a4paper]{article}
\usepackage[margin=24mm]{geometry}
\usepackage{amsmath,amssymb,bm}
\usepackage{microtype}
\usepackage[hidelinks]{hyperref}
\title{Resonant Chemistry\\Shell-Level N-Body Topology with Nuclear Centres\\Candidate Formalism v0.1}
\author{Adrian Lipa}
\date{13 August 2026}
\begin{document}
\maketitle
\begin{abstract}
We formalize a shell-level extension of many-electron dynamics in which nuclei remain explicit attractive Coulomb centres and topology is introduced as a diagnostic of trajectory/history structure rather than a new force. The construction separates three-dimensional exchange permutations from projection-dependent winding and reserves braid-group language for effective two-dimensional or constrained sectors. Particle-hole duality is defined for arbitrary $s,p,d,f$ subshells and linked to a falsifiable conformal reduction with anchor $z_\star=1/2$.
\end{abstract}
\section{Epistemic status}
Established control physics, representation definitions, implemented diagnostics, validated results, and promoted canon are kept distinct. This document is a candidate formalism and is not canon.
\section{Control system}
For nuclei $A=1,\dots,M$ and electrons $i=1,\dots,N$,
\begin{equation}
H=\sum_i\frac{\mathbf p_i^2}{2}-\sum_{i,A}\frac{Z_A}{|\mathbf r_i-\mathbf R_A|}+\sum_{i<j}\frac{1}{|\mathbf r_i-\mathbf r_j|}+\cdots.
\end{equation}
The topology layer does not alter this Hamiltonian by default.
\section{Nuclear centres}
Each nucleus $\mathcal A_A=(Z_A,\mathbf R_A)$ defines $V_A=-Z_A/|\mathbf r-\mathbf R_A|$. In conservative dynamics this is an attractive potential centre, not a dissipative attractor. ``Attractor'' is used literally only for reduced SCF/relaxation maps with basins.
\section{Shell duality}
For orbital angular momentum $l$, capacity is $C_l=2(2l+1)$. Occupation $k$ has hole partner $k^\star=C_l-k$ and the half-filled self-dual point is $k=2l+1$. Hence $p^1\leftrightarrow p^5$, $p^2\leftrightarrow p^4$, $p^3\leftrightarrow p^3$; similarly $d^5$ and $f^7$ are self-dual.
\section{Topology boundary}
In three spatial dimensions the exchange topology of identical point particles is represented by permutations. Braid-group invariants are introduced only in an effective two-dimensional or constrained sector. For a declared projection $P:\mathbb R^3\to\mathbb C$ we may define
\begin{equation}
w_{iA}=\frac{1}{2\pi}\oint d\arg P(\mathbf r_i-\mathbf R_A),\qquad
w_{ij}=\frac{1}{2\pi}\oint d\arg P(\mathbf r_i-\mathbf r_j).
\end{equation}
These are explicitly projection-dependent.
\section{Conformal reduction}
Given a declared shell-history coordinate $\Phi[\gamma]\in\mathbb C$,
\begin{equation}
\zeta=\Phi[\gamma]-z_\star,\qquad \eta=\frac{1}{\zeta}.
\end{equation}
The candidate anchor $z_\star=1/2$ must survive replacement-anchor and Möbius controls before physical interpretation.
\section{Shell-topology record}
A minimal diagnostic state is
\begin{equation}
\Sigma_{nlk}=\left(n,l,k,C_l-k,\pi_{\rm exch},\{w_{iA}\},\{w_{ij}\},\Gamma_{\rm geom},\eta\right).
\end{equation}
\section{Validation ladder}
The required order is: (1) recover SCF/CI controls; (2) extract topology diagnostics without fitting experimental energies; (3) test $p^k\leftrightarrow p^{6-k}$; (4) test transfer $2p^k\to3p^k$; (5) only then compare frozen shell classes with knot/conformal families.
\section{Implementation}
The accompanying module \texttt{reschem/shell\_nbody\_topology.py} implements shell capacity, particle-hole duality, nuclear Coulomb centres, exchange parity, declared-plane nuclear and pair winding diagnostics, and the conformal inverse coordinate.
\section{Current verdict}
\textbf{IMPLEMENTED candidate representation; local unit tests PASS; physical validation OPEN; canonical promotion DENIED pending shell-only tests.}
\end{document}
